\subsection{An explicit rational eight-cubic seed}
\label{app:eight-cubic-smooth-seed}

\begin{theorem}
\label{thm:eight-cubic-smooth-seed}
Over $\mathbb Q$, and over every sufficiently large prime field, there is a word on $16$ distinct coordinates whose maximum agreement with cubics is $7$ and whose complete nearest list consists of exactly eight cubics. Over arbitrarily large prime fields there is also a length-$32$, dimension-$7$ Reed--Solomon word whose maximum agreement is $14$ and whose complete nearest list consists of exactly eight codewords. This agreement exceeds the audited first-order curve at rate $7/32$. The construction has fixed length and list size; no growing-list or asymptotic-tightness assertion follows.
\end{theorem}

\begin{proof}
We first give the rational source as sections of $\mathcal O_{\mathbb P^1}(3)$, then specify an ordinary affine Reed--Solomon chart. Number the coordinates from $0$ to $15$ and put
\[
 (z_j)_{j=0}^{15}=\left(
 \frac{11}{9},\frac{11}{5},\frac{77}{65},1,
 \frac{66}{65},\frac{22}{25},\frac65,\frac{55}{52},
 \frac{11}{10},\frac{77}{95},\infty,0,
 \frac{33}{5},\frac{33}{25},\frac{99}{95},\frac{33}{35}\right).
\]
The eight sections $f_i(X)=\sum_{r=0}^3 c_{ir}X^r$ and their exact agreement sets are
\[
\begin{array}{c|rrrr|l}
 i&c_{i0}&c_{i1}&c_{i2}&c_{i3}& S_i\\\hline
0&693/25&-165/2&883/11&-6175/242&2,4,5,9,13,14,15\\
1&-99/14&90/7&-81/14&0&0,3,7,10,12,14,15\\
2&-198/25&69/5&-65/11&0&4,6,8,10,12,13,14\\
3&0&0&0&0&0,3,4,9,10,11,13\\
4&0&-1&1&0&2,3,5,8,10,11,14\\
5&0&-12/5&38/11&-130/121&2,4,7,8,11,12,15\\
6&0&6&-54/11&0&0,1,2,6,10,11,12\\
7&0&-6&116/11&-50/11&1,3,6,8,11,13,15
\end{array}
\]
At infinity, a section is evaluated by its $X^3$ coefficient. Define $v_j=f_i(z_j)$ for any $i$ with $j\in S_i$. Direct rational arithmetic verifies that this definition is independent of the choice of $i$, and that $f_i$ agrees precisely on $S_i$.

To obtain sixteen finite nodes, use
\[
 u_j=\frac1{z_j-3},\qquad
 q_i(U)=U^3 f_i(3+U^{-1}).
\]
The convention is $u_{10}=0$. Set $w_j=u_j^3v_j$ for finite $z_j$, and $w_{10}=v_{10}=0$. These expressions define rational polynomials and values, and preserve all incidences. The nodes are distinct and remain distinct modulo $17$.

Here is a finite certificate of completeness. In the residue field let
\[
 R(Y)=3+Y+15Y^2+7Y^3,\qquad
 P_i(Y)=10^{-1}q_i((11-Y)/3)+R(Y).
\]
The nodes $11-3u_j$ reduce to $j+1$, and the corresponding received values reduce, after this same gauge change, to
\[
 (9,3,12,15,0,14,8,13,11,13,11,8,10,14,10,13).
\]
Enumeration of all $\binom{16}{4}=1820$ four-point cubic interpolants gives $1022$ distinct polynomials. The maximum agreement is seven, and exactly the eight $P_i$ attain it. This also proves completeness over the algebraic closure of the residue field: four matching points determine a cubic over $\mathbb F_{17}$.

Any characteristic-zero cubic with at least seven agreements on the rational source is determined by four of these rational points and is integral at $17$, since their Vandermonde determinant is a unit. Its reduction is a known seven-match cubic. Its matching support must therefore be exactly that known support, and four-point uniqueness identifies it with the corresponding $q_i$. This proves the complete rational source claim. Since there are only finitely many four-point interpolants and nonzero evaluation differences, the same complete list persists modulo every sufficiently large prime.

For the square pullback it is essential to use the affine cover
\[
 U=\frac{11-T^2}{3},\qquad
 T=\pm\sqrt{11-3u_j},
\]
rather than $U=T^2$. The sixteen radicands reduce to $1,\ldots,16$ modulo $17$. Thus all thirty-two nodes have distinct reductions in $\mathbb F_{289}$, and the eight rational sextics $q_i((11-T^2)/3)$ reduce to
\[
 10\bigl(P_i(T^2)-R(T^2)\bigr).
\]
After the common polynomial subtraction and nonzero scalar change, this is exactly the certified square pullback of the residue seed.

For clarity, its completeness certificate is as follows. A non-even degree-six polynomial has the form $E(Y)+T O(Y)$ with $Y=T^2$, $\deg E\le3$, $\deg O\le2$, and $O\ne0$. At most two fibers can contribute both signs, so fourteen matches require at least twelve distinct hit fibers. At a hit fiber,
\[
 w^2+B(Y)w+C(Y)=0,\qquad
 B=-2E,\quad C=E^2-YO^2.
\]
The eleven coefficients of $B,C$ satisfy twelve linear equations. Among all $\binom{16}{12}=1820$ choices, $1686$ systems have rank eleven and are inconsistent, $132$ have rank eleven and are consistent, and two have rank ten and are inconsistent. There are $84$ distinct consistent solutions. Every one fails at least one necessary identity $C_0=E_0^2$ or $C_6=E_3^2$. Hence no non-even competitor exists, even over the algebraic closure. Even competitors are covered by the cubic census. The complete residue nearest list therefore consists of eight sextics, each with fourteen matches.

For a characteristic-zero competitor, interpolation on seven matching nodes makes its coefficients integral at the chosen unramified quadratic place above $17$. Its reduction is one of those eight sextics with exactly fourteen matches. Seven-point uniqueness then identifies the competitor with the corresponding rational sextic. This argument also handles a nonzero odd part whose reduction vanishes. Thus the square pullback has a complete eight-element nearest list and maximum agreement fourteen.

Finally adjoin the finitely many square roots to a number field and exclude the finitely many primes where a coefficient, node difference, or nonzero evaluation difference of a seven-point interpolant degenerates. Arbitrarily large rational primes split completely in its Galois closure, yielding the claimed ordinary prime-field codes.
\end{proof}

The first-order expression at $\rho=7/32$, $a=14/32$ is
\[
 (8-\rho)a^2-6\rho a+\rho(4\rho-5)=\frac{105}{8192}>0.
\]
The exact rational table and all $56$ incidence identities are checked by
\nolinkurl{research/relaxed_eight_cubic_route/rational_seed.py}.
The coefficientwise residue gauge and exact cover are independently recorded by
\texttt{verify\_rational\_cover.py}. The complete source and pullback enumerations have independent interpolation and elimination replays in
\texttt{independent\_verify.py} and \texttt{verify\_quadratic\_pullback.py}.
